% d5_exam.tex -- DISTILLED: exam prep. Every practice-exam and Sai-midterm question in one line,
% recurring question patterns with answer recipes + tiny worked numbers, 60-second facts, HW1-3 answers.
% (Full question text / official solutions: d6_practice_exam.tex and d7_sai_midterm.tex; full HW: raw_hw.tex.)
\section{Exam prep: past questions, patterns, quick answers}
\ixh{model evaluation}\ixh{model comparison}\ixh{exam patterns}

{\small Format (Canvas): $\le$12 multiple choice + 4 short answer, paper notes only; R code / R output may appear.
Sai's midterm (Summer 2026, 26 pts) = 7 MC $\times$ 2 pts + 2 free-response (8 + 4 pts). MC rule from the cover sheet:
\emph{more than one bubble = zero}; ``Both X and Y'' options are traps unless \emph{both} parts are right. Round to 3--4 decimals.\par}

%======================================================================
{\small Every practice-exam and Sai question, in full with worked answers: the next two sections.\par}

\subsection{Question patterns and how to answer them}
\ixh{exam patterns}

\begin{topic}{P1. Confusion matrix $\to$ metrics \src{PE Q17; HW3 Q10; S3.34--37}}
\ixh{confusion matrix}\ixh{true positive}\ixh{false positive}
\textbf{Recipe.} (1) Find the orientation: are \emph{rows} predicted or true? (PE Q17: rows = predicted;
S3.34: rows = true; \ixc{conf\_mat()} prints \texttt{Prediction} rows $\times$ \texttt{Truth} columns.)
(2) Which class is ``positive''? (yardstick: the \emph{first} factor level.) (3) Label TP, FP, FN, TN by
\emph{meaning}. (4) Divide by the right total: actual positives (TPR, FNR), actual negatives (FPR, spec.),
predicted positives (precision).
\par\textbf{Example.} Truth Yes: 40 predicted Yes, 10 predicted No; truth No: 20 predicted Yes, 130 No. So TP 40, FN 10,
FP 20, TN 130, $n=200$.
\ix{accuracy} $=170/200=0.85$; \ix{sensitivity} $=40/50=0.80$; \ix{specificity} $=130/150=0.867$;
\ix{FPR} $=20/150=0.133$; \ix{precision} $=40/60=0.667$; error rate $=0.15$.
\end{topic}

\begin{topic}{P2. Change the threshold: what happens? \src{Sai Q7--8; PE Q15}}
\ixh{threshold}
\textbf{Recipe.} Predict positive if $\hat p\ge t$. Raise $t$ $\Rightarrow$ fewer positives $\Rightarrow$ TP, FP
$\downarrow$; \ix{sensitivity}/TPR $\downarrow$, FPR $\downarrow$, \ix{specificity} $\uparrow$; the
ROC curve and AUC are \textbf{unchanged} (only the operating point slides down-left). Lower $t$: the reverse.
Coefficients and $\hat p$ never change with $t$.
\par\textbf{Example.} $\hat p$ / truth: 0.15 No, 0.40 Yes, 0.62 No, 0.85 Yes.\\
{\footnotesize\begin{tabular}{@{}l c c c c c c@{}}
$t$ & TP & FP & FN & TN & sens & spec\\ \midrule
0.3 & 2 & 1 & 0 & 1 & 1.0 & 0.5\\
0.5 & 1 & 1 & 1 & 1 & 0.5 & 0.5\\
0.7 & 1 & 0 & 1 & 2 & 0.5 & 1.0\\
\end{tabular}}
\end{topic}

\begin{topic}{P3. Read an ROC curve / compare AUCs \src{PE Q18--19; Sai Q3, Q8c; ROC}}
\ixh{ROC curve}\ixh{AUC}
\textbf{Recipe.} $x$ = FPR ($1-$spec.), $y$ = TPR (sens.); each point = one threshold. Up/right = lower $t$.
Pick the point by costs: FN costly (disease, fraud) $\Rightarrow$ high-TPR point (up/right); FP costly $\Rightarrow$
low-FPR point (down/left). Compare models by AUC (higher better; 0.5 = coin flip, 1 = perfect, AUC $=1$ $\Rightarrow$
suspect \ix{target leakage}). AUC does not depend on the chosen threshold.
\par\textbf{Example.} Same 4 cases as P2: positives score 0.40, 0.85; negatives 0.15, 0.62. Of the $2\times2$
positive--negative pairs, the positive scores higher in 3 $\Rightarrow$ \ix{AUC} $=3/4=0.75$ (standard ``probability a
random positive outranks a random negative'' reading). Model A AUC 0.88 vs.\ B 0.79 $\Rightarrow$ A
discriminates better at almost every threshold.
\end{topic}

\begin{topic}{P4. Flexibility, bias--variance, $k$ in KNN \src{PE Q1--2; Sai Q1--2, Q5; S1.113--115}}
\ixh{bias-variance tradeoff}\ixh{flexibility}
\textbf{Recipe.} More flexible $\Rightarrow$ lower \ix{bias}, higher \ix{variance}, training error $\downarrow$, test
error U-shaped. \emph{More} flexible: smaller $k$, add predictors / polynomial / interaction terms, QDA instead of
LDA, KNN instead of a linear model. Overfit symptom: train $\gg$ test $\Rightarrow$ make it \emph{less} flexible
(increase $k$, drop terms).
\par\textbf{Example.} $k=1$: each training point is its own neighbor, training error $=0$ (maximal variance).
$k=n$: always predicts the overall majority class (maximal bias). With 400 training points, $k=1$
train/test error 0.00/0.30 vs.\ $k=15$ 0.18/0.21 $\Rightarrow$ choose $k=15$.
\end{topic}

\begin{topic}{P5. Pick a model from a train/test table \src{PE Q3; HW2 Q9; Lab 2}}
\ixh{training error}\ixh{testing error}
\textbf{Recipe.} Choose by the \textbf{test} column (highest accuracy/AUC/$R^2$, lowest RMSE/MAE). Then diagnose:
big train--test gap = \ix{overfitting}; train $\approx$ test but both poor = \ix{underfitting}.
\par\textbf{Example.} RMSE train/test: linear 2.18/1.95, KNN 1.46/2.21 (HW2 abalone) $\Rightarrow$ linear model
(lower test RMSE); KNN overfits.
\end{topic}

\begin{topic}{P6. KNN by hand \src{PE Q5--6; S1.103--109}}
\ixh{KNN}\ixh{Euclidean distance}
\textbf{Recipe.} Standardize first; compute distances $d=\sqrt{\sum_j(x_{j}-x_{0j})^2}$; sort; classification =
majority vote of the $k$ closest; regression = \emph{average} $y$ of the $k$ closest.
\par\textbf{Example.} $x_0=(0,0)$; $(1,1)$ A: $d=1.41$; $(2,0)$ B: $2$; $(0,3)$ B: $3$; $(3,4)$ A: $5$.
$k=1$ $\Rightarrow$ A; $k=3$ $\Rightarrow$ B (A,B,B). Regression: nearest 3 have $y=10,12,17$ $\Rightarrow\hat y=13$.
\end{topic}

\begin{topic}{P7. Interpret a logistic coefficient / compute $\hat p$ \src{PE Q14, Q16; Sai Q7; S3.12--25}}
\ixh{logistic regression}\ixh{odds ratio}\ixh{log odds}
\textbf{Recipe.} $\beta_j$ = change in \ix{log odds} per unit $X_j$, others held fixed; $e^{\beta_j}$ = \ix{odds ratio}
(multiplicative); only the \emph{sign} transfers to $p$; say ``associated with'', never ``causes''. To predict:
$z=\beta_0+\sum\beta_jx_j$, $\hat p=e^z/(1+e^z)$, then compare with $t$.
\par\textbf{Example.} $\beta_0=-3$, $\beta_1=0.05$ (age). Age 40: $z=-1$, $\hat p=e^{-1}/(1+e^{-1})=0.269$
$\Rightarrow$ ``No'' at $t=0.5$. Age $+10$: odds $\times e^{0.5}=1.65$. $\beta=0.7$: odds $\times e^{0.7}=2.01$ (about double).
Template: ``A one-unit increase in $X$ is associated with a $\beta$ increase in the log odds of $Y$=Yes
(odds multiplied by $e^{\beta}$), holding other predictors fixed.''
\end{topic}

\begin{topic}{P8. Read regression output; dummy and interaction coefficients \src{HW2 Q6; PE Q13; S2}}
\ixh{linear regression}\ixh{dummy variable}
\textbf{Recipe.} \ixc{tidy()} columns: \texttt{term}, \texttt{estimate} ($\hat\beta$), \texttt{std.error},
\texttt{statistic} $=$ estimate/SE (\ix{t-statistic}), \texttt{p.value} ($<0.05$ $\Rightarrow$ significant).
Dummy coefficient = difference from the \ix{reference level} (the level with no column), others fixed.
Interaction $X_1{:}X_2$ coefficient = how much the $X_1$ slope changes per unit of $X_2$.
\par\textbf{Example.} \texttt{diameter} 2.20, SE 0.308 $\Rightarrow t=7.14$, $p\approx10^{-12}$. Salary $=40+2\,\text{exp}+5\,
\text{IT}-3\,\text{HR}$ (Sales = reference): IT with 10 yrs $\Rightarrow 40+20+5=65$; HR earns 8 less than IT at equal
experience.
\end{topic}

\begin{topic}{P9. Spot an interaction in a table / plot \src{PE Q9; Sai Q6}}
\ixh{interaction term}
\textbf{Recipe.} Compute the effect of $X_1$ at each level of $X_2$ (differences, or slopes). Different $\Rightarrow$
interaction; equal (parallel lines) $\Rightarrow$ additive. Correlation between predictors is irrelevant.
\par\textbf{Example.} Low/High dose: A 20/30, B 25/50 $\Rightarrow$ +10 vs +25: interaction. A 20/30, B 25/35
$\Rightarrow$ +10 both: none (B just shifted up 5).
\end{topic}

\begin{topic}{P10. Which variables may be predictors? \src{PE Q7, Q10, Q12; Sai Q4; HW2 Q3}}
\ixh{target leakage}\ixh{identifier variable}
\textbf{Recipe} --- ask four questions of each variable: (1) Is it the \ix{outcome} or computed from it?
(2) Is it known \emph{at prediction time}? (no $\Rightarrow$ \ix{target leakage}). (3) Is it unique per row
(\ix{identifier variable})? (4) Is it an exact function of another predictor (\ix{collinearity})?
Group labels with a few shared levels (hospital, dorm, major) are fine $\Rightarrow$ dummy-encode.
\par\textbf{Example.} Predict abalone \texttt{age}: \texttt{rings} is out (age $=$ rings $+1.5$, leakage).
Loan default: \texttt{default\_date} out (only exists after default), \texttt{application\_id} out (ID).
\end{topic}

\begin{topic}{P11. Choose recipe steps (and their order) \src{S2; Lab 2--3; HW2 Q3; HW3 Q4}}
\ixch{recipes}\ixh{feature engineering}
{\footnotesize
\begin{tabularx}{\linewidth}{@{}>{\raggedright\arraybackslash\leavevmode}X >{\raggedright\arraybackslash\leavevmode}X@{}}
\toprule
Situation & Step\\ \midrule
missing numeric values & \ixc{step\_impute\_linear()} (\ixc{imp\_vars()}) or \ixc{step\_impute\_mean()}\\
categorical predictor, lm/glm/LDA & \ixc{step\_dummy()} (default $K-1$ columns)\\
categorical, KNN & \texttt{step\_dummy(..., one\_hot = TRUE)} OK (no intercept)\\
KNN (distance) & \ixc{step\_normalize()} = \ixc{step\_center()} + \ixc{step\_scale()}\\
effect of $X_1$ depends on $X_2$ & \ixc{step\_interact()}, after dummying: \texttt{\textasciitilde{} starts\_with("sex"):fare}\\
curved relationship & \ixc{step\_poly()}; skewed $X$ \ixc{step\_log()}\\ \bottomrule
\end{tabularx}}
\par Order: impute $\to$ dummy $\to$ interact $\to$ normalize. Recipe is \emph{trained on the training data}
(\texttt{recipe(y \textasciitilde{} ., data = train)}).
\end{topic}

\begin{topic}{P12. Spot the bug in tidymodels code \src{PE Q7, Q13; Sai Q9; Lab 3}}
\ixch{workflow()}
Checklist: \textbf{(1)} recipe built but model fit with \texttt{fit(y \textasciitilde{} ., data)} $\Rightarrow$ recipe ignored;
use \ixc{workflow()} \pipe{} \ixc{add\_recipe()} \pipe{} \ixc{add\_model()}. \textbf{(2)} \texttt{y \textasciitilde{} .} includes IDs /
leakage. \textbf{(3)} KNN without \ixc{step\_normalize()}. \textbf{(4)} \texttt{one\_hot = TRUE} with an intercept model
(rank deficient). \textbf{(5)} \ixc{nearest\_neighbor()} without \ixc{set\_mode()} (mode unknown). \textbf{(6)} metrics
on \emph{training} data / \ixc{predict()} on training instead of \texttt{testing()}. \textbf{(7)} \ixc{roc\_auc()}
given \ixc{.pred\_class} instead of the probability column (\ixc{.pred\_Yes}, from \texttt{type = "prob"}).
\textbf{(8)} wrong event level (first factor level is the event). \textbf{(9)} no \ixc{set.seed()} before
\ixc{initial\_split()}; no \ixc{strata} for an imbalanced outcome. \textbf{(10)} \ixc{step\_interact()} naming a
variable that \ixc{step\_dummy()} already renamed (\texttt{sex} $\to$ \texttt{sex\_male}).
\end{topic}

\begin{topic}{P13. LDA vs.\ QDA vs.\ logistic vs.\ KNN \src{PE Q20; Sai Q2, Q5; S3.89--93}}
\ixh{linear discriminant analysis (LDA)}\ixh{quadratic discriminant analysis (QDA)}
\textbf{Recipe.} Boundary linear? $\to$ logistic or LDA (LDA if predictors roughly normal within classes / small $n$;
logistic if normality doubtful). Class covariances clearly differ \emph{and} $n$ large relative to $p$ $\to$ QDA.
Covariances equal $\to$ LDA (QDA only adds variance). Very irregular boundary, large $n$, small $p$ $\to$ KNN.
Flexibility: KNN $>$ QDA $>$ LDA $=$ logistic.
\par\textbf{Example.} $n=150$, $p=8$, similar spreads $\Rightarrow$ LDA. $n=50{,}000$, $p=2$, different
covariances $\Rightarrow$ QDA. Binary outcome, skewed/dummy predictors $\Rightarrow$ logistic regression.
\end{topic}

\begin{topic}{P14. Correlation, collinearity, class imbalance \src{PE Q4, Q8; Sai Q9; HW3 Q1--2}}
\ixh{Pearson correlation}\ixh{class imbalance}
\textbf{Recipe.} $r$ measures \emph{linear} association only ($r\approx0$ with a U-shape is possible). Two predictors
with $|r|\approx1$ (or one an exact rescaling of the other) $\Rightarrow$ \ix{collinearity}: keep one. Imbalanced
outcome $\Rightarrow$ stratify the split (\ixc{strata}) and judge accuracy against the majority-class baseline.
\par\textbf{Example.} 90\% negatives: ``always negative'' scores \ix{accuracy} 0.90 with sensitivity 0 and AUC 0.5.
Titanic: 61.6\% died $\Rightarrow$ any useful model must beat 0.616.
\end{topic}

\begin{topic}{P15. Read other R output \src{Lab 2--3; HW2--3}}
\ixh{R output}
\ixc{predict()} default \texttt{type = "class"} $\to$ \ixc{.pred\_class}; \texttt{type = "prob"} $\to$
\ixc{.pred\_Yes}/\texttt{.pred\_No}; regression $\to$ \texttt{.pred}. \ixc{augment()} = data + all prediction columns.
Metric tibbles: \texttt{.metric}, \texttt{.estimator}, \ixc{.estimate}. LDA print: \emph{Prior probabilities} = training
class proportions ($\hat\pi_k$), \emph{Group means} = $\hat\mu_k$, \emph{LD1} = discriminant direction. kknn print:
``Minimal misclassification'' = leave-one-out training error. \ixc{initial\_split()} print
\texttt{<Training/Testing/Total>}.
\end{topic}

%======================================================================
\subsection{60-second facts}
\begin{keybox}[High-yield one-liners]
\begin{itemize}[leftmargin=1em,itemsep=0pt]
\item Expected test MSE\ixh{mean squared error (MSE)}\ixh{bias-variance tradeoff}{} $=\mathrm{Var}(\hat f)+\mathrm{Bias}^2+\mathrm{Var}(\epsilon)$; last term = \ix{irreducible error}, always $\ge$ it.
\item \ix{supervised learning} has a $Y$; \ix{unsupervised learning} does not. \ix{regression} = numeric $Y$, \ix{classification} = categorical $Y$.
\item Parametric (linear, logistic, LDA, QDA) = assume a form for $f$; \ix{non-parametric model} (KNN) = no form, needs more data.
\item \ix{KNN}: small $k$ = flexible\ixh{flexibility}{} = low bias / high variance; $k=1$ training error 0. Always normalize\ixh{normalization}.
\item \ix{curse of dimensionality}: with many predictors, the ``nearest'' neighbors are far away $\Rightarrow$ KNN degrades.
\item Lower training error is never by itself a reason to prefer a model; compare on \ix{testing data}.
\item \ix{Bayes classifier}: assign the class with the largest $\Pr(Y=k\mid X=x_0)$; its error = \ix{Bayes error rate} (lowest possible).
\item $R^2$\ixh{R-squared}{} = fraction of variance in $Y$ explained; RMSE/MAE\ixh{RMSE}\ixh{MAE}{} in units of $Y$, lower is better.
\item $K$ levels $\Rightarrow$ $K-1$ dummies\ixh{dummy variable}\ixh{one-hot encoding}{} with an intercept; $K$ (one-hot) only without an intercept / for KNN.
\item Interaction\ixh{interaction term}{} = effect of one predictor depends on another; \ix{hierarchical principle}: keep main effects if the interaction is in.
\item Correlation $\ne$ interaction; correlation $\approx 0$ $\ne$ useless; correlation $\ne$ causation.
\item \ix{logistic regression}: \ix{logit} linear in $X$; $e^{\beta}$ = \ix{odds ratio}; $p=0.5\Leftrightarrow$ log odds $=0$; odds $=p/(1-p)$.
\item Logistic fit by \ix{maximum likelihood}; LDA by plugging in $\hat\pi_k,\hat\mu_k,\hat\Sigma$.
\item \ix{LDA}: common $\Sigma$\ixh{covariance matrix}, linear \ix{decision boundary}; \ix{QDA}: $\Sigma_k$ per class, quadratic boundary, more parameters.
\item \ix{confounding}: a coefficient's sign can flip once another predictor is added (Default \texttt{student}).
\item \ix{sensitivity} $=$ \ix{TPR} $=$ recall $=TP/(TP+FN)$; \ix{specificity} $=TN/(TN+FP)=1-$\ix{FPR}; \ix{precision} $=TP/(TP+FP)$.
\item Raise \ix{threshold}: sens $\downarrow$, spec $\uparrow$, positives $\downarrow$; ROC curve and AUC unchanged.
\item \ix{ROC curve}: $x$ = FPR, $y$ = TPR; diagonal = random (AUC 0.5); top-left corner = perfect.
\item \ix{AUC} is threshold-free and summarizes all thresholds; AUC = 1 on real data $\Rightarrow$ check for \ix{target leakage}.
\item \ixc{yardstick}: first factor level\ixh{event level}\ixch{factor()}{} = event $\Rightarrow$ \texttt{factor(survived, levels = c("Yes","No"))}.
\item \ixc{glm} (logistic) models the \emph{second} level: with Yes first, coefficients are for ``No''.
\item Stratify\ixh{stratified sampling}{} on the outcome so train/test keep the class proportions; \ixc{set.seed()} for \ix{reproducibility}.
\item Recipe + model $\to$ \ixc{workflow()} $\to$ \ixc{fit()} on training $\to$ \ixc{augment()}/\ixc{predict()} on testing $\to$ metrics.
\item \ixc{step\_normalize()} cannot fix collinearity; it only puts predictors on a common scale.
\item A predicted probability\ixh{posterior probability}{} is \emph{not} an accuracy; a model has accuracy, one patient has a probability.
\end{itemize}
\end{keybox}

%======================================================================
\subsection{HW answers at a glance}
\ixh{homework}
\begin{topic}{HW1 --- ML ideas; EDA on \texttt{mpg} \src{HW1; submission}}
\ixch{mpg}
\textbf{Q1} supervised\ixh{supervised learning}{} = labeled $Y$ to learn $f$; unsupervised\ixh{unsupervised learning}{} = no $Y$, find structure (clusters).
\textbf{Q2} \ix{regression}: quantitative $Y$; \ix{classification}: qualitative $Y$ (class or class probabilities).
\textbf{Q3} regression: \ix{RMSE}, $R^2$\ixh{R-squared}{} (\ix{MAE}); classification: \ix{accuracy}, \ix{sensitivity}/\ix{recall} (AUC).
\textbf{Q4} \ix{descriptive model}: visually emphasize a trend; \ix{inferential model}: which features are significant,
test theories/relationships; \ix{predictive model}: best prediction of $Y$, minimal reducible error \src{S1.50}.
\textbf{Q5} \ix{mechanistic model} = \ix{parametric model} (assumed form; easier to understand; higher \ix{bias}, lower \ix{variance});
\ix{empirically-driven model} = \ix{non-parametric model} (few assumptions, more data, more flexible: lower bias, higher variance); both
estimate $f$ and both can overfit\ixh{overfitting}. (Correction to my submission: empirical models are \emph{lower}-bias, not ``unbiased''.)
\textbf{Q6} (a) chance a voter supports the candidate = \textbf{predictive}\ixh{prediction}; (b) effect of personal contact = \textbf{inferential}\ixh{inference}.
\textbf{EDA:} \ixc{hwy} \ix{histogram} bimodal (peaks $\approx$16 and $\approx$26), slightly right-skewed\ixh{skewness}; \texttt{hwy}--\texttt{cty}
\ix{scatterplot} strong positive linear ($r=0.96$); \ix{bar plot}: most cars dodge (37), fewest lincoln (3); \ix{boxplot}: \texttt{hwy} falls with \texttt{cyl}
(medians 29/24/17 for 4/6/8); \ixc{corrplot()} \ix{correlation matrix}: \texttt{displ}--\texttt{cyl} 0.93, mpg vs.\ \texttt{displ}/\texttt{cyl}
$\approx -0.8$, \texttt{year} $\approx 0$.
\end{topic}

\begin{topic}{HW2 --- abalone, linear regression vs.\ KNN \src{HW2; submission}}
\ixch{abalone}
\textbf{Q1} \texttt{age = rings + 1.5} (\ixc{mutate()}); mean 11.4, median 10.5, SD 3.22, range 2.5--30.5: unimodal, right-skewed,
$\approx$6.7\% boxplot outliers\ixh{outlier}. \textbf{Q2} \ixc{initial\_split()} \texttt{(prop = 0.85, strata = age)}\ixh{stratified sampling}\ixch{prop}\ixch{strata}{} $\to$ 3548 / 629.
\textbf{Q3} recipe: \ixc{step\_dummy()} $\to$ \ixc{step\_interact()} (type$\times$shucked, longest$\times$diameter,
shucked$\times$shell) $\to$ \ixc{step\_center()} $\to$ \ixc{step\_scale()}; exclude \texttt{rings}: age is an exact
function of it (\ix{target leakage}). \textbf{Q4--6} \ixc{linear\_reg()} \texttt{"lm"}, \ixc{nearest\_neighbor()} \texttt{(neighbors = 7)}
\texttt{"kknn"}\ixch{kknn}, two workflows\ixch{workflow()}, fit on training. \textbf{Q7} hypothetical female (\ixc{predict()}): $\hat y=21.8$ years.
\textbf{Q8} \ixc{metric\_set()}\texttt{(rsq, rmse, mae)} + \ixc{augment()} on test: LM $R^2$ 0.585, RMSE 1.95, MAE 1.44; KNN 0.483 / 2.21 / 1.62. $R^2=0.585$: 58.5\% of the variance
in test-set age explained. \textbf{Q9} LM wins on test; KNN training $R^2$ 0.811 $\to$ test 0.483 = \ix{overfitting}
(flexible, high variance); LM 0.552 $\to$ 0.585 (rigid, stable). \textbf{Q10--12 (231)} \ix{reducible error} $=\mathrm{Var}+\mathrm{Bias}^2$,
\ix{irreducible error} $=\mathrm{Var}(\epsilon)$; both nonnegative $\Rightarrow$ test error $\ge \mathrm{Var}(\epsilon)$.
\end{topic}

\begin{topic}{HW3 --- Titanic classifiers (MY OWN answers; not submitted/graded) \src{HW3; raw HW section}}
\ixch{titanic}
{\color{warnc}\itshape Self-worked solutions (R, seed 3435); exact numbers depend on the seed/split --- know the reasoning.}
\textbf{Q1} 70/30, \ixc{strata} \texttt{= survived}: 623 / 268. \ix{missing data}: \texttt{age} 128 (impute), \texttt{cabin} 481
(drop), \texttt{embarked} 2; IDs\ixh{identifier variable}{} (\texttt{passenger\_id}, \texttt{name}, \texttt{ticket}) not predictors.
Stratify\ixh{stratified sampling}{} because the outcome is imbalanced\ixh{class imbalance}{} (38.4\% Yes) --- both sets keep the same proportion.
\textbf{Q2} sex is a strong predictor (survival 72.9\% female vs.\ 20.3\% male); class too (63.0 / 47.7 / 23.8\% for
1st/2nd/3rd). Percent stacked bars\ixh{percent stacked bar chart}{} (\ixc{geom\_bar()}, \texttt{position = "fill"}) compare \emph{proportions} when the bar totals differ.
\textbf{Q3} \texttt{sib\_sp}--\texttt{parch} $+0.43$, \texttt{age}--\texttt{sib\_sp} $-0.33$, \texttt{fare} weakly $+$
with the family counts; all $|r|<0.45$\ixh{Pearson correlation}, no \ix{collinearity} problem (needs \texttt{use = "pairwise.complete.obs"}).
\textbf{Q4} \ixc{step\_impute\_linear()}\ixh{linear imputation}\texttt{(age, impute\_with = imp\_vars(...))}\ixch{imp\_vars()}{} $\to$ \texttt{step\_dummy(all\_nominal\_predictors())}\ixch{step\_dummy()}\ixch{all\_nominal\_predictors()}{}
$\to$ \texttt{step\_interact(\textasciitilde{} starts\_with("sex"):fare + age:fare)}\ixch{step\_interact()}\ixch{starts\_with()}\ixh{interaction term}.
\textbf{Q5--8} \ixc{logistic\_reg()} \texttt{"glm"}\ixch{glm}, \ixc{discrim\_linear()} / \ixc{discrim\_quad()} \texttt{"MASS"}\ixch{MASS}{}
(\ixc{discrim}), \texttt{nearest\_neighbor(neighbors = 10)}\ixch{neighbors}{} \texttt{"kknn"}; same recipe, four workflows.
\textbf{Q9} \ixc{predict()} \texttt{type = "prob"} + \ixc{bind\_cols()} $\to$ \ixc{roc\_auc()}; training \ix{AUC}: logistic 0.845, LDA 0.844, QDA 0.852, KNN 0.964 (inflated: own neighbor).
\textbf{Q10} testing AUC\ixh{testing data}: logistic 0.880, \textbf{LDA 0.881}, QDA 0.859, KNN 0.879 $\Rightarrow$ LDA (near-tie).
LDA test \ix{confusion matrix} (\ixc{conf\_mat()}) TP 73, FP 22, FN 30, TN 143: accuracy 0.806, \ix{sensitivity} 0.709, \ix{specificity} 0.867; \ixc{roc\_curve()} + \ixc{autoplot()} for the plot.
Train 0.844 vs.\ test 0.881: rigid low-variance model, the difference is split-to-split sampling noise (small test
set); KNN's 0.964 $\to$ 0.879 is the overfitting pattern. \textbf{Q11--12 (231)} \ix{logit} $=$ inverse \ix{logistic function};
$x_1+2$ multiplies \ix{odds} by $e^{2\beta_1}$ (\ix{odds ratio}); $\beta_1<0$: $p\to0$ as $x_1\to\infty$, $p\to1$ as $x_1\to-\infty$.
\end{topic}
